% GENERATED FILE. Edit canonical lessons, then run npm run generate:books.
% canonical-source-hash: fnv1a64:dcd1789f7fe9fa38
% canonical-lessons: RU-C180-golovnaya-bol, RU-C180-son, RU-C180-diyeta, RU-C180-vitamin, RU-C180-lyogkiye

\chapter{Headache, Sleep, Diet, Vitamin, Lungs}
\label{ch:ru-headache-sleep-diet-vitamin-lungs}
\hlchaptermodality{180}

\begin{chapteropening}
\textbf{By the end of this chapter:} \emph{I can say a headache, sleep, a diet, a vitamin and the lungs.}

Complete the last lesson of chapter 180: I can say a headache, sleep, a diet, a vitamin and the lungs.
\end{chapteropening}

\section[\texorpdfstring{\ru{головная} \ru{боль} (golovnáya bol)}{головная боль (golovnáya bol)}]{\ru{головная} \ru{боль} (golovnáya bol) — a headache}

\label{lesson:RU-C180-golovnaya-bol}



\begin{quote}
Before the new one: say the Russian for a bruise, then the Russian for a fracture.
\end{quote}

\subsection*{You'll want to know: \ru{головная} \ru{боль}}
\textbf{\ru{головная} \ru{боль}} (\emph{golovnáya bol}) — \textquotedblleft{}a headache\textquotedblright{}.

\begin{cousinweb}
\textbf{\ru{головная} \ru{боль}} is \textquotedblleft{}head pain\textquotedblright{}: \emph{\ru{головная}}, from \emph{\ru{голова}}, \textquotedblleft{}head\textquotedblright{}, with \emph{\ru{боль}}, \textquotedblleft{}pain\textquotedblright{}. Vasmer relates \emph{\ru{боль}} to Old English \emph{bealu}, \textquotedblleft{}evil\textquotedblright{}, a link Berneker rejected.
\end{cousinweb}

\begin{grammarlens}[title={What you've built}]
One more word for this chapter.
\end{grammarlens}

\subsection*{Your turn}
\begin{itemize}
\raggedright
  \item \emph{Say it:} \emph{golovnáya bol}
  \item \emph{Say it:} \emph{golovnáya bol}, once more
  \item \emph{Say it:} say \emph{golovnáya bol}
  \item \emph{Recall:} say the Russian for a bruise, then the Russian for a fracture, then say \emph{golovnáya bol} again
\end{itemize}

\begin{tcolorbox}[breakable,colback=teal!4,colframe=teal!35!black,title={Before you move on}]
What does \ru{головная} \ru{боль} mean? (\textquotedblleft{}A headache\textquotedblright{}.) Say it once more.
\end{tcolorbox}
\section[\texorpdfstring{\ru{сон} (son)}{сон (son)}]{\ru{сон} (son) — sleep; a dream}

\label{lesson:RU-C180-son}



\begin{quote}
Before the new one: say the Russian for a fracture, then the Russian for a headache.
\end{quote}

\subsection*{You'll want to know: \ru{сон}}
\textbf{\ru{сон}} (\emph{son}) — \textquotedblleft{}sleep; a dream\textquotedblright{}.

\begin{cousinweb}
\textbf{\ru{сон}} $\leftarrow$ Proto-Slavic \emph{*sŭnŭ}, from \emph{*sŭpnŭ}, with \emph{\ru{спать}}, \textquotedblleft{}to sleep\textquotedblright{} — inherited. Vasmer lists Lithuanian \emph{sãpnas}, \textquotedblleft{}a dream\textquotedblright{}, Sanskrit \emph{svápna-}, Greek \emph{hýpnos} and Latin \emph{somnus}, \textquotedblleft{}sleep\textquotedblright{}.
\end{cousinweb}

\begin{grammarlens}[title={What you've built}]
One more word for this chapter.
\end{grammarlens}

\subsection*{Your turn}
\begin{itemize}
\raggedright
  \item \emph{Say it:} \emph{son}
  \item \emph{Say it:} \emph{son}, once more
  \item \emph{Say it:} say \emph{son}
  \item \emph{Recall:} say the Russian for a fracture, then the Russian for a headache, then say \emph{son} again
\end{itemize}

\begin{tcolorbox}[breakable,colback=teal!4,colframe=teal!35!black,title={Before you move on}]
What does \ru{сон} mean? (\textquotedblleft{}Sleep; a dream\textquotedblright{}.) Say it once more.
\end{tcolorbox}
\section[\texorpdfstring{\ru{диета} (diyéta)}{диета (diyéta)}]{\ru{диета} (diyéta) — a diet}

\label{lesson:RU-C180-diyeta}



\begin{quote}
Before the new one: say the Russian for a headache, then the Russian for sleep; a dream.
\end{quote}

\subsection*{You'll want to know: \ru{диета}}
\textbf{\ru{диета}} (\emph{diyéta}) — \textquotedblleft{}a diet\textquotedblright{}.

\begin{cousinweb}
\textbf{\ru{диета}} is a bookish loan from French \emph{diète} or Latin \emph{diaeta}, from Greek \emph{díaita}, \textquotedblleft{}a way of living\textquotedblright{} (Vasmer).
\end{cousinweb}

\begin{grammarlens}[title={What you've built}]
One more word for this chapter.
\end{grammarlens}

\subsection*{Your turn}
\begin{itemize}
\raggedright
  \item \emph{Say it:} \emph{diyéta}
  \item \emph{Say it:} \emph{diyéta}, once more
  \item \emph{Say it:} say \emph{diyéta}
  \item \emph{Recall:} say the Russian for a headache, then the Russian for sleep; a dream, then say \emph{diyéta} again
\end{itemize}

\begin{tcolorbox}[breakable,colback=teal!4,colframe=teal!35!black,title={Before you move on}]
What does \ru{диета} mean? (\textquotedblleft{}A diet\textquotedblright{}.) Say it once more.
\end{tcolorbox}
\section[\texorpdfstring{\ru{витамин} (vitamín)}{витамин (vitamín)}]{\ru{витамин} (vitamín) — a vitamin}

\label{lesson:RU-C180-vitamin}



\begin{quote}
Before the new one: say the Russian for sleep; a dream, then the Russian for a diet.
\end{quote}

\subsection*{You'll want to know: \ru{витамин}}
\textbf{\ru{витамин}} (\emph{vitamín}) — \textquotedblleft{}a vitamin\textquotedblright{}.

\begin{cousinweb}
\textbf{\ru{витамин}} is an international coinage of 1912: Latin \emph{vīta}, \textquotedblleft{}life\textquotedblright{}, with \emph{amine}, a chemical name.
\end{cousinweb}

\begin{grammarlens}[title={What you've built}]
One more word for this chapter.
\end{grammarlens}

\subsection*{Your turn}
\begin{itemize}
\raggedright
  \item \emph{Say it:} \emph{vitamín}
  \item \emph{Say it:} \emph{vitamín}, once more
  \item \emph{Say it:} say \emph{vitamín}
  \item \emph{Recall:} say the Russian for sleep; a dream, then the Russian for a diet, then say \emph{vitamín} again
\end{itemize}

\begin{tcolorbox}[breakable,colback=teal!4,colframe=teal!35!black,title={Before you move on}]
What does \ru{витамин} mean? (\textquotedblleft{}A vitamin\textquotedblright{}.) Say it once more.
\end{tcolorbox}
\section[\texorpdfstring{\ru{лёгкие} (lyógkiye)}{лёгкие (lyógkiye)}]{\ru{лёгкие} (lyógkiye) — the lungs}

\label{lesson:RU-C180-lyogkiye}



\begin{quote}
Before the new one: say the Russian for a diet, then the Russian for a vitamin.
\end{quote}

\subsection*{You'll want to know: \ru{лёгкие}}
\textbf{\ru{лёгкие}} (\emph{lyógkiye}) — \textquotedblleft{}the lungs\textquotedblright{}.

\begin{cousinweb}
\textbf{\ru{лёгкие}} is \textquotedblleft{}the light ones\textquotedblright{}, from \emph{\ru{лёгкий}}, \textquotedblleft{}light\textquotedblright{}: the lungs are the lightest organ, as in English \emph{lights}, \textquotedblleft{}lungs\textquotedblright{}.
\end{cousinweb}

\begin{grammarlens}[title={What you've built}]
One more word for this chapter.
\end{grammarlens}

\subsection*{Your turn}
\begin{itemize}
\raggedright
  \item \emph{Say it:} \emph{lyógkiye}
  \item \emph{Say it:} \emph{lyógkiye}, once more
  \item \emph{Say it:} say \emph{lyógkiye}
  \item \emph{Recall:} say the Russian for a diet, then the Russian for a vitamin, then say \emph{lyógkiye} again
\end{itemize}

\begin{tcolorbox}[breakable,colback=teal!4,colframe=teal!35!black,title={Before you move on}]
What does \ru{лёгкие} mean? (\textquotedblleft{}The lungs\textquotedblright{}.) Say it once more.
\end{tcolorbox}
